% !TeX program = pdflatex
\documentclass[10pt,a4paper]{article}

\usepackage[utf8]{inputenc}
\usepackage[T1]{fontenc}
\usepackage{geometry}
\usepackage[hidelinks]{hyperref}
\usepackage{xcolor}

\geometry{
  top=1.7cm,
  bottom=1.8cm,
  left=1.8cm,
  right=1.8cm
}

\pagestyle{empty}
\frenchspacing
\setlength{\parindent}{0pt}
\setlength{\parskip}{0.45em}

\definecolor{mydarkgreen}{RGB}{0,0,0}

\begin{document}

\begin{minipage}[t]{0.48\textwidth}
\textbf{\large Javier Andres TARAZONA JIMENEZ}\\
Étudiant ingénieur IA (Bac +5) / Data Science\\
Massy, Île-de-France, France\\
\href{mailto:javitar06@gmail.com}{javitar06@gmail.com}\\
+33 07 46 36 97 75\\
\href{https://www.linkedin.com/in/javier-andres-tarazona-jimenez-84b489222/}{LinkedIn}
\end{minipage}
\hfill
\begin{minipage}[t]{0.42\textwidth}
\raggedleft
Société Générale Assurances\\
DataLab\\
Courbevoie, Île-de-France, France\\
\end{minipage}

\vspace{1.1cm}

\raggedleft
Massy, le 19 juin 2026

\vspace{0.7cm}

\raggedright
\textbf{Objet : Candidature à l'alternance Data Scientist IA Générative / Computer Vision}

\vspace{0.5cm}

{\small

Madame, Monsieur,

Actuellement étudiant colombien en double diplôme entre l’Universidad Nacional de Colombia et ENSTA Paris, au sein de l’Institut Polytechnique de Paris, je souhaite vous présenter ma candidature pour l’alternance Data Scientist IA Générative / Computer Vision au sein du DataLab de Société Générale Assurances, à partir de septembre 2026.

J’ai choisi l’informatique par volonté de comprendre des systèmes complexes, de construire des solutions concrètes à partir d’idées abstraites, et de développer une autonomie technique solide dans un cadre académique exigeant. L’intelligence artificielle s’est ensuite imposée comme une continuité naturelle : passer de la programmation de systèmes à la conception de systèmes capables de percevoir, apprendre et aider à la décision. Ce domaine m’attire particulièrement parce qu’il confronte l’ingénieur à des problèmes ouverts, ambigus et fortement liés au réel. La vision par ordinateur, en particulier, correspond à mon intérêt pour les applications concrètes : transformer des mathématiques, des données et du code en perception du monde. Plus largement, l’IA représente pour moi une technologie transversale, à la fois scientifique et pragmatique, capable d’avoir un impact visible sur la société et l’économie.

Ce qui m’intéresse chez Société Générale Assurances n’est pas seulement la dimension d’un grand acteur financier, mais surtout la combinaison de quatre éléments essentiels : l’échelle, la donnée, la régulation et l’impact métier. Votre activité ne se limite pas aux assurances : elle touche aux décisions financières, au risque, au comportement client, à la distribution digitale, à la protection patrimoniale et à la relation de confiance avec les assurés. Cet environnement constitue, à mes yeux, un cadre particulièrement stimulant pour l’IA, car il transforme des problèmes abstraits en décisions concrètes de business.

L’IA y est utile non comme effet de mode, mais comme outil appliqué à des enjeux réels : classification documentaire, extraction d’informations, personnalisation, automatisation, chatbots, détection de signaux récurrents, expérience client, fiabilité opérationnelle et aide à la décision. Contrairement à un environnement purement technologique ou à une startup centrée uniquement sur l’expérimentation, Société Générale Assurances impose un niveau d’exigence supérieur : les modèles doivent être robustes, explicables, auditables, industrialisables et compatibles avec un contexte réglementé. C’est précisément cette contrainte qui rend le poste attractif pour moi.

Mon profil est en forte adéquation avec vos besoins. Lors de mon stage de recherche à ENSTA Paris – U2IS, je travaille sur le pré-entraînement auto-supervisé avec PyTorch et Vision Transformers, la génération de données 3D contrôlées avec Python et Blender, le suivi de configurations expérimentales et l’évaluation de représentations. Chez APEDYS91, j’ai développé une application NLP responsable combinant règles, LLMs locaux, spaCy/NER et contrôles sur les entités et les nombres. Chez Engin A.I, j’ai renforcé mes compétences en vision par ordinateur, analyse de données et qualité logicielle Python, avec des pipelines OpenCV/NumPy, du clustering non supervisé et une refactorisation orientée maintenabilité.

Je peux ainsi apporter à vos équipes une base technique solide en Python, PyTorch, NLP, computer vision, APIs, Docker, Git et développement logiciel, mais aussi une capacité à relier expérimentation, rigueur méthodologique et besoins métiers. Mon parcours international en Colombie, au Canada et en France m’a également appris à travailler dans des contextes multiculturels, avec autonomie, curiosité et sens du résultat.

Rejoindre Société Générale Assurances représenterait pour moi une étape structurante : passer d’une IA principalement académique et expérimentale à une IA industrialisée, responsable et utile à grande échelle. Je serais honoré de pouvoir contribuer à vos projets et d’apprendre auprès de vos data scientists, ML engineers et experts métiers.

Je vous prie d’agréer, Madame, Monsieur, l’expression de mes salutations distinguées.

\vspace{0.8cm}

\textbf{Javier Andres TARAZONA JIMENEZ}

}

\end{document}